\section{Introduction}\label{sec:intro}

\subsection{The theorems}

The Lean file \leanfile{Main} proves five statements. In ordinary notation they are the following.

\begin{theorem}[\lean{LFunction_ne_zero_of_half_lt_re}]\label{intro:thm:grh}
Let $N \ge 1$, let $\chi$ be a Dirichlet character modulo $N$, and let $s \in \C$ with
$\Re s > 1/2$. If $\chi$ is not principal or $s \ne 1$, then $L(s,\chi) \ne 0$.
\end{theorem}

\begin{theorem}[\lean{LFunction_eq_zero_re}, \lean{LFunction_eq_zero_re_of_isPrimitive}]
\label{intro:thm:line}
Let $\chi$ be a Dirichlet character modulo $N \ge 1$, primitive or not. Every zero $s$ of
$L(s,\chi)$ with $0 < \Re s < 1$ has $\Re s = 1/2$.
\end{theorem}

\begin{theorem}[\lean{riemannHypothesis}, \lean{riemannZeta_eq_zero_re}]\label{intro:thm:rh}
Every zero $s \ne 1$ of the Riemann zeta function other than $-2, -4, -6, \dots$ has real part
$1/2$.
\end{theorem}

The statements use Mathlib's definitions of $\zeta$ and of Dirichlet $L$-functions;
\lean{riemannHypothesis} has Mathlib's own type \lean{RiemannHypothesis}, and the other four have
the types fixed in \texttt{grh/src/Statements.lean}. Theorems~\ref{intro:thm:line}
and~\ref{intro:thm:rh} follow from Theorem~\ref{intro:thm:grh} by the functional equation
(Proposition~\ref{zeros:prop:line}). Theorems~\ref{intro:thm:grh} and~\ref{intro:thm:line} say
nothing about zeros of $L(s,\chi)$ with $\Re s \le 0$. In Theorem~\ref{intro:thm:rh} the point
$s = 1$ is excluded because Mathlib's $\zeta$ has a value there.

\subsection{The source, and what this document is}

The proof exists as Lean~4 code built on Mathlib: the repository \texttt{lsernau/grh-lean}, folder
\texttt{grh}, the ``readable edition'', with 803 Lean files in \texttt{grh/src} and about
$300{,}000$ lines of code. Its \texttt{README} states that the proof is Levent Alp\"oge's, with the
AI system Claude; that the Lean text and the accompanying documents were written by Claude; and
that this edition is a shorter rewriting of Alp\"oge's Lean proof that follows its method
(it refers to \url{https://alpo.ge/remark.pdf}). The code compiles without \texttt{sorry}, and the
five theorems depend only on the axioms \texttt{propext}, \texttt{Classical.choice} and
\texttt{Quot.sound}. The repository reports that an export of the five theorems with all they
depend on is accepted by three independent checkers (nanoda, con-leche and con-ron). The
\texttt{README} at the top of the repository says: ``No mathematician has read it through, and
no paper comes with it.''

This document is a reading of that code. Its purpose is to say, in ordinary mathematical
language, what the argument is. The reader should know exactly what it is and what it is not.

\begin{itemize}
\item \emph{It follows the Lean text.} Every definition and every stated result is tied to a
  Lean declaration. In the text the declaration appears as a short marker; for instance
  \lean{New.star_inequality} is the main inequality, \texttt{New.star\_inequality}. The index in
  Appendix~\ref{sec:app-leanindex} gives the name behind each marker, and in the PDF each marker
  links to the declaration in the repository, at the file and line. Files are cited like
  \leanfile{New/Reduction/MainInequality}, relative to the folder \texttt{grh/src}. Where this
  document and the Lean text differ, the Lean text is right.
\item \emph{It admits the background.} The code is divided into a folder \texttt{New}, the
  argument, and a folder \texttt{Lib}, results classed as known in the literature that Mathlib
  lacks. We write out \texttt{New}. The results of \texttt{Lib} on which the argument turns are
  stated as \emph{Facts}, without proof, with the sources that the Lean files name; other
  declarations of \texttt{Lib} are cited where a proof uses them. Appendix~\ref{sec:app-facts}
  lists both kinds, and it is complete: it is generated from the source of this document. The
  code often proves a weaker or differently shaped statement than the textbook theorem; we state
  what the code proves. We have not compared the Facts with their sources. The list of
  references at the end gives the sources that the Lean files name.
\item \emph{It is not a second verification.} The correctness of the theorems rests on the Lean
  kernel, not on this text. The statements here were translated from the formal statements by
  reading, and the proofs are explanations of the formal proofs: the mechanism, the key identity
  or inequality, and how the pieces combine. A step that is routine or long is replaced by a
  reference to the Lean declaration that proves it. A translation can contain mistakes that the
  kernel cannot catch.
\item \emph{How it was made.} The repository credits the Lean text to Claude, the author of
  the present text. The two are separate pieces of work: this text was written from the
  published repository alone, with no access to the work that produced the code. It was drafted
  from the Lean sources:
  the documentation in the files, the formal statements and, where needed, the proofs. Each
  section was drafted from one group of folders. Each was then checked, independently of its
  drafting, against the Lean text: every stated definition and result was compared with the
  formal declaration it cites, and the proofs of the main results with the structure of the
  formal proofs. That check examined about $680$ statements, definitions and proofs and
  corrected about $250$ of them; most corrections restored a hypothesis that the draft had
  dropped, and some replaced a proof step that described a mechanism other than the one the
  code uses. Part~1 and the four sections that carry most of the proof
  (Sections~\ref{sec:der}, \ref{sec:cg-selmer}, \ref{sec:cg-first} and~\ref{sec:cg-lifting})
  were then checked a second time, with the Lean proofs read for the steps that the first check
  had compared only as statements; this corrected about $50$ further items and added the
  explanation of about $50$ steps. A third reader re-derived Part~1 and the arithmetic of the
  budgets of Parts~2 to~4 by hand and compared about $45$ statements with their declarations;
  that found one discrepancy. Last, the notation was unified and repeated definitions were
  replaced by references, in a pass that was to change no mathematics and was itself compared
  with the text before it. None of this shows that no error remains. The comments in the code
  were trusted unless the code contradicted them. All citations of Lean names were checked
  mechanically to name declarations that exist.
\item \emph{What is not explained.} For many auxiliary steps the text gives the statement and
  the Lean name but not the reason, and for some it describes the mechanism only as far as the
  documentation in the code does. Appendix~\ref{sec:app-open} lists these steps section by
  section.
\end{itemize}

\subsection{The proof in outline}

Figure~\ref{fig:overview} shows the top of the proof. Each of Parts~2, 3 and~4 opens with a
diagram of the same kind for the theorem it proves
(Figures~\ref{fig:part2}, \ref{fig:part3}, \ref{fig:part4a} and~\ref{fig:part4b}).

% Diagram of the top of the proof (Introduction): from the final theorems down to Theorems A, B, C.
\begin{figure}[tbp]
\centering
\begin{tikzpicture}[x=1cm,y=1cm]
% row 0: the statements on the critical line
\node[pfmain] (line) at (0,0) {Thms.~\ref{intro:thm:line}, \ref{intro:thm:rh}\\zeros on the critical line};
% row 1
\node[pfmain] (grh) at (0,-1.7) {Thm.~\ref{main:thm:grh}\\(=~Thm.~\ref{intro:thm:grh})\\half-plane theorem};
\node[pfnode] (fe) at (6,-1.7) {Prop.~\ref{zeros:prop:line}\\functional equation};
% row 2
\node[pfnode] (crit) at (-3.75,-3.4) {Thm.~\ref{zeros:thm:criterion}\\one-sided criterion};
\node[pfnode] (one) at (3.75,-3.4) {Thm.~\ref{main:thm:onesided}\\one-sided bound};
% row 3
\node[pffact, text width=1.9cm] (landau) at (-6.25,-5.1) {Fact~\ref{zeros:fact:landau}\\Landau's lemma};
\node[pffact, text width=1.9cm] (prime) at (-3.75,-5.1) {Fact~\ref{zeros:fact:prime}\\choice of the prime};
\node[pfnode, text width=1.9cm] (poles) at (-1.25,-5.1) {Prop.~\ref{zeros:prop:poles}\\poles};
\node[pfnode, text width=1.9cm] (atx) at (1.25,-5.1) {Prop.~\ref{main:prop:atx}\\remainder at the first scale};
\node[pfnode, text width=1.9cm] (exist) at (3.75,-5.1) {Thm.~\ref{setup:thm:existence}\\existence of the data};
\node[pfnode, text width=1.9cm] (pnt) at (6.25,-5.1) {Thm.~\ref{zeros:thm:pnt}\\prime number theorem};
% row 4
\node[pfnode] (rem) at (-1.75,-6.8) {Prop.~\ref{setup:prop:rem}\\total remainder};
\node[pfmain] (star) at (1.25,-6.8) {Thm.~\ref{main:thm:star}\\main inequality};
% row 5
\node[pfmain] (A) at (-1.75,-8.5) {Thm.~\ref{main:thm:A}\\Theorem A: upper bound};
\node[pfmain] (B) at (1.25,-8.5) {Thm.~\ref{main:thm:B}\\Theorem B: spanning};
\node[pfmain] (C) at (4.25,-8.5) {Thm.~\ref{main:thm:C}\\Theorem C: lower bound};
% row 6: the parts
\node[pfpart] (P2) at (-1.75,-10.2) {Part~2\\seesaw identity,\\divisibility, size\\Fig.~\ref{fig:part2}};
\node[pfpart] (P3) at (1.25,-10.2) {Part~3\\second moment\\Fig.~\ref{fig:part3}};
\node[pfpart] (P4) at (4.25,-10.2) {Part~4\\chain of lifts,\\twists\\Figs.~\ref{fig:part4a}, \ref{fig:part4b}};
% arrows: premise -- conclusion
\draw[pfedge] (grh) -- (line);
\draw[pfedge] (fe) -- (line);
\draw[pfedge] (crit) -- (grh);
\draw[pfedge] (one) -- (grh);
\draw[pfedge] (landau) -- (crit);
\draw[pfedge] (prime) -- (crit);
\draw[pfedge] (poles) -- (crit);
\draw[pfedge] (atx) -- (one);
\draw[pfedge] (exist) -- (one);
\draw[pfedge] (pnt) -- (one);
\draw[pfedge] (rem) -- (atx);
\draw[pfedge] (star) -- (atx);
\draw[pfedge] (A) -- (star);
\draw[pfedge] (B) -- (star);
\draw[pfedge] (C) -- (star);
\draw[pfedge] (P2) -- (A);
\draw[pfedge] (P3) -- (B);
\draw[pfedge] (P4) -- (C);
\end{tikzpicture}
\caption{The top of the proof: from the theorems of the introduction down to Theorems~A, B
and~C (Sections~\ref{sec:intro}, \ref{sec:zeros}, \ref{sec:setup} and~\ref{sec:main}). Each box
is a result, with its number and a short name; a box with several numbers stands for a group of
results. An arrow from one box to another means that the proof of the second, as written in this
paper, uses the first. Bold boxes are the final theorems,
the main inequality and Theorems~A, B and~C. Dashed grey boxes are background results admitted
without proof. Double boxes stand for the parts that prove Theorems~A, B and~C; each has its own
diagram.}
\label{fig:overview}
\end{figure}


\subsubsection*{A one-sided bound suffices}
Fix a modulus $q_0$ and a prime $d$, and let $K = \Q(\sqrt{-d})$ with quadratic character
$\chiK$. Consider the weighted sum over the prime powers $n \equiv 1 \bmod q_0$ with
$\chiK(n) \ne 1$,
\[
  R_{q_0}(x) = \sum_{\substack{n \le x \\ n \equiv 1 \bmod q_0}}
     \Lambda(n)\bigl(1 - \chiK(n)\bigr)\Bigl(\frac1n - \frac1x\Bigr) - \frac{\log x}{\varphi(q_0)} - c ,
\]
with the constant $c$ for which $R_{q_0}(x) \to 0$. A theorem of Landau's type
(Section~\ref{sec:zeros}) shows that an \emph{upper} bound $R_{q_0}(x) \le x^{-1/2+\eps}$ for all
large $x$, for every $\eps > 0$ and all primes $d$ in a residue class, excludes zeros of the
$L$-functions modulo $q_0$ in $\Re s > 1/2$. The rest of the proof establishes this upper bound for one $x$ at a time.

\subsubsection*{Size against divisibility}
The upper bound is proved by the principle that drives transcendence theory and Stepanov's proof
of the Riemann hypothesis for curves: a non-zero integer cannot be both small and highly
divisible. One builds a lattice of ``integral'' objects, shows that it contains many vectors that
are small, and compares the resulting upper bound for a determinant with a lower bound that comes
from the product formula. In his remark, Alp\"oge describes the argument as a descendant of
Stepanov's, in which the high multiplicity of an auxiliary function is replaced by a chain of
theta lifts, each lift contributing one unit of multiplicity, and the length of the chain is about
$\sqrt{x}$.

The primes enter through divisibility. A theta function of an elliptic curve with complex
multiplication by $K$ has Taylor coefficients (``jets'') at the origin which, suitably
normalised, are algebraic integers; and at a prime $p$ that does \emph{not} split in $K$ a jet of
order $a$ is divisible by about $p^{a\gamma_p}$, with $\gamma_p = 1/p + 1/p^3 + \cdots$ for inert
$p$ and $\gamma_d = \frac12(1/d + 1/d^2 + \cdots)$ (Section~\ref{sec:der}). The reason is a
Frobenius congruence $\theta(\pi_p z)/\theta(z)^{q_p} \equiv 1$ modulo the prime $\pi_p$ of $K$
above $p$, of norm $q_p$: after a scaling of $z$ it makes the theta series integral up to an
exponential factor, and Legendre's formula for $a!$ does the rest. Summed over $p$ with weights
$\log p$ these exponents give half of $\sum_n \Lambda(n)(1 - \chiK(n))\cdot a/n$, a sum over
primes of the kind that occurs in $R_{q_0}$ (the factor $2$ returns because sizes are squared
norms).

\subsubsection*{The lattice}
For each large integer $x$ the proof fixes \emph{admissible data} (Section~\ref{sec:setup}): a
cyclotomic extension $E = K(\zeta)$, a CM field with $f$ infinite places; hermitian spaces over
$E$; and scales $x < x_1 < b$ with $a = xb$. \emph{Source forms} are algebraic modular forms of
weight about $a$ on a definite unitary group $\U(X_0)$. \emph{Target forms} are automorphic forms
on a unitary group $\U(Y_0)$. A Poincar\'e series construction maps source forms to target forms,
and a chain of $m \approx \sqrt{x}$ theta lifts $\Theta_\vartheta$, twisted by a character
$\vartheta$, maps target forms to algebraic modular forms on a definite group $\U(Y_m)$, where
integrality and size make sense. The target forms with the maps $\Theta_\vartheta$ are the
lattice; a family of target forms has a \emph{covolume} (its size) and an \emph{index} (its
distance from integrality).

\subsubsection*{Three theorems}
\emph{Theorem A} (the upper bound, Part~2) produces short integral vectors. A seesaw identity
rewrites the lifts of a Poincar\'e series as a combination of chains of lifts between definite
spaces, whose values are jets of CM theta functions. Their archimedean size is explicit, and their
divisibility at the non-split primes is as above. Dividing by the guaranteed power of each prime
leaves forms whose lifts have logarithmic size at most $\hmain - m\,\rem + f a\log^4 b$, where
the \emph{main amount} $\hmain$ is explicit and huge and the \emph{total remainder} $\rem$ is a combination of the remainders
$R_{q_0}$.

\emph{Theorem B} (spanning, Part~3) shows by a second-moment computation, of the type of the
Petersson formula, that few translates of the Poincar\'e map span almost all target forms.

\emph{Theorem C} (the lower bound, Part~4) bounds covolume plus twice the index from below by
$\dT\,\hmain$, $\dT$ the dimension of the target, up to small errors, for every independent family
of almost full rank. It evaluates a Gram determinant by the product formula, with a different
twist at each prime (Theorem~\ref{setup:thm:twist}). The Gram determinant $Z(\vartheta)$ of a
chain of lifts is not evaluated: exchange identities of Siegel--Weil type write $Z(1)/Z(\vartheta)$,
up to a power of the norm of the modulus, as ratios of twisted to untwisted Hecke $L$-values times
the same ratios for the $m$ single lifts; the latter are bounded through comparison operators and
lifts of characters of $\U(1)$, whose squared norms are Hecke $L$-values (Rallis inner product
formula). The main amount reappears as the height of $Z(1)/Z(\vartheta)$, with $\vartheta$ chosen
for each prime; the twists are chosen through congruences, Galois representations and Selmer
groups, which bound the denominators of the lifts.

\subsubsection*{The conclusion}
In the comparison the main amount cancels, and what remains is $m\cdot\rem \le f a\log^5 b$
(Theorem~\ref{main:thm:star}). The weights of the forms are chosen larger at a sparse set of
places of $E$ of class $-1$ modulo $q_0$; this is how the residue class modulo $q_0$ enters, and
it makes $\rem$ equal to $a\,\#\Phi_b\,R_{q_0}(x)$ up to remainders at the far larger point $x_1$,
where the classical prime number theorem suffices (Proposition~\ref{setup:prop:rem}). Dividing by
$m \approx \sqrt{x}$ gives $R_{q_0}(x) \le C x^{-1/2+6\eps}$ for data with $\log b \approx x^\eps$.

\subsection{Organisation}

Part~1 is the reduction: the analytic criterion (Section~\ref{sec:zeros}), the data and the
general lower bound for Gram determinants (Section~\ref{sec:setup}), and the derivation of the
theorem from Theorems A, B and C (Section~\ref{sec:main}); it is complete in itself, given
Theorems B and C and the results of Part~2 from which Section~\ref{sec:main} assembles Theorem~A. Parts~2, 3 and~4 prove Theorems A, B and C. Table~\ref{intro:tab:parts} lists the
sections with the folders of the Lean text that they cover. The appendices give the notation
(Appendix~\ref{sec:app-notation}), the background results admitted
(Appendix~\ref{sec:app-facts}), the steps cited without explanation
(Appendix~\ref{sec:app-open}) and the index of Lean declarations
(Appendix~\ref{sec:app-leanindex}); the references follow.

The parts are of very unequal weight. A few sections quote a percentage for a result: it is the
share, among the declarations that the final theorem \lean{riemannHypothesis} needs, of those
that are needed only through that result. The figures come from the dependency graph of the
compiled code, in which the final theorem needs $49{,}815$ declarations. By this measure almost
half ($46\%$) of the proof is needed only through Theorem~C, and most of that only through the
existence of suitable twists (Section~\ref{sec:cg-selmer}).

\begin{table}[ht]
\centering
\small
\begin{tabular}{llp{7.9cm}}
\toprule
section & subject & Lean folders under \texttt{New/} \\
\midrule
\ref{sec:zeros} & one-sided criterion & \texttt{Reduction/Zeros}, \texttt{Main} \\
\ref{sec:setup}, \ref{sec:data} & data, heights, lattices & \texttt{Reduction/Data}, \texttt{Reduction/Lattices} \\
\ref{sec:main} & main inequality & \texttt{Reduction} \\
\midrule
\ref{sec:ub-objects} & objects & \texttt{UpperBound/Kernels}, \texttt{UpperBound/Objects} \\
\ref{sec:der} & derivatives of theta functions & \texttt{Derivatives} \\
\ref{sec:ub-seesaw} & seesaw identity & \texttt{UpperBound/Seesaw} \\
\ref{sec:ub-divis} & divisibility and size & \texttt{UpperBound/Divisibility}, \texttt{UpperBound/Size} \\
\ref{sec:ub-eldiv} & elementary divisors & \texttt{UpperBound/ElementaryDivisors} \\
\midrule
\ref{sec:sp-moment} & second moment, Bessel functions & \texttt{Spanning/SecondMoment}, \texttt{Multiplicities}, \texttt{Bessel}; \texttt{Appendix/DiscreteSeries} \\
\ref{sec:sp-unfold} & unfolding & \texttt{Spanning/Unfolding} \\
\ref{sec:sp-main} & main term & \texttt{Spanning/MainTerm} \\
\ref{sec:sp-frames} & local estimates & \texttt{Spanning/Frames} \\
\ref{sec:sp-remain} & remaining terms & \texttt{Spanning/RemainingTerms} \\
\midrule
\ref{sec:det-chain} & Theorem C, chain of lifts & \texttt{Determinant/Twist}, \texttt{Norms}, \texttt{Chain}, \texttt{AuxiliaryForms} \\
\ref{sec:det-exchange} & exchange identity & \texttt{Determinant/Exchange} \\
\ref{sec:ol-values} & one lift, Hecke values & \texttt{OneLift} except \texttt{ComparisonOperator} \\
\ref{sec:ol-compare} & comparison operator & \texttt{OneLift/ComparisonOperator} \\
\ref{sec:cg-selmer} & twists, Selmer groups & \texttt{Congruences/FitTwistsExist}, \texttt{Selmer}, \texttt{GaloisRepresentations} \\
\ref{sec:cg-first} & Ribet's method & \texttt{Congruences/FirstOrder} \\
\ref{sec:cg-lifting} & local loss & \texttt{Congruences/LocalLoss}, \texttt{Lifting} \\
\bottomrule
\end{tabular}
\caption{The sections and the Lean folders they cover.}
\label{intro:tab:parts}
\end{table}

\subsection{Conventions}

\emph{Total functions.} In Lean every function is defined everywhere. Mathlib's $\zeta$ has a
value at $s = 1$, its $\Gamma$ has the value $0$ at the poles, division by zero gives $0$. Formal
statements therefore carry side conditions such as $s \ne 1$; we keep them when they matter and
otherwise write as usual.

\emph{Places.} $\Aut(\C)$ is the group of all field automorphisms of $\C$. A pair $\iota$ of a
prime $\ell$ and a field isomorphism $\C \isom \Qlbar$ is written $\iota \mid \ell$. A complex
number is integral at $\iota$ if its image has absolute value at most $1$; a quantity is bounded
``at every conjugate'' if the bound holds after applying every $\sigma \in \Aut(\C)$.

\emph{Letters.} The symbols listed in Appendix~\ref{sec:app-notation} have one meaning in the
whole paper; among them are the letters of the data of Section~\ref{sec:setup}. There are two
standing exceptions: $s$ is also a complex variable (Section~\ref{sec:zeros}) and an element of
$\Gbar$ (as in $s\nPoin$), and a bound variable may reuse a letter. Sections introduce further
symbols for their own objects and say what they mean; where a section gives a listed letter a
local meaning, it says so at that place. A twist is $\vartheta$. Constants are explicit in the
Lean text and are not optimised; we keep them.

\emph{Sizes.} Inequalities between the sizes of the data are used constantly. They compare
quantities through bounds for their logarithms, which reduces them to linear inequalities between
$\log x$, $\log x_1$, $\log b$, $\log\log b$, $\log\log x_1$ and a constant; they are proved by one
uniform mechanism (\lean{New.Data.le_of_hasLogBound}), and are quoted without proof.
